\uuid{mb3j}
\titre{Distance}
\auteur{bodin ; hanani ; mzari}
\organisation{exo7}
\licence{CC BY-NC-SA 4.0}
\datecreate{2020-01-20}
\niveau{L1}
\module{Géométrie}
\chapitre{Géométrie de l'espace}
\sousChapitre{Distance}
\difficulte{1}
\theme{Distance}
\exo7qcmid{lille-1.geometrie-espace.021}

\contenu{
\texte{
Pour ces questions, l'espace est muni d'un repère orthonormé direct  $(O, \vec{i}, \vec{j}, \vec{k})$.
}

\begin{question}
\qclassification{141.05}
\qtags{motcle=distance}
Soit $A(1,1,1)$ et $P$ le plan d'équation cartésienne : $x+y+z+1=0$. La distance de $A$ à $P$ est :
\begin{answers}

    \bad {$\frac{1}{\sqrt 3}$}

    \bad {$\frac{2}{\sqrt 3}$}

   \bad{$\sqrt 3$}

    \good{$\frac{4}{\sqrt 3}$}

\end{answers}
\begin{explanations}
Si $P$ est d'équation  $ax+by+cz+d=0$ et $A(x_0,y_0,z_0)$, la distance de $A$ à $P$ est donnée par : $\frac{|ax_0+by_0+cz_0+d|}{\sqrt{a^2+b^2+c^2}}$.
\end{explanations}
\end{question}
}
